\documentclass[10pt,a4paper]{amsart}
\usepackage[margin=19mm]{geometry}
\usepackage{amsmath,amssymb,mathtools}
\usepackage[scr=rsfs,cal=euler]{mathalfa}
\usepackage{xfrac}
\usepackage[unicode,hidelinks,bookmarks=false]{hyperref}
\newcommand{\ie}{i.e.\ }
\newcommand{\cf}{cf.\ }
\newcommand{\resp}{respectively\ }
\newcommand{\Subsection}{\subsection}
\providecommand{\nobreakdash}{\nobreak\hbox{-}}
\setlength{\emergencystretch}{2em}
\multlinegap0pt
\usepackage{amssymb}
\usepackage{mathtools}
\newtheorem{lemma}{Lemma}
\newtheorem{theorem}{Theorem}
\newcommand{\R}{{\mathbb{R}}}
\title{Detrended Fluctuation Analysis for Continuous Real Variable Functions}
\author{Luis Gil-Maqueda, Benjam\'in A. Itz\'a-Ortiz}
\date{Source: arXiv:2203.15940v3 (CC BY 4.0)}
\begin{document}
\maketitle

\noindent\emph{Source and licence.} Detrended Fluctuation Analysis for Continuous Real Variable Functions. Version arXiv:2203.15940v3 (CC BY 4.0), distributed by arXiv under the Creative Commons Attribution 4.0 International licence, http://creativecommons.org/licenses/by/4.0/, as recorded in the arXiv licence metadata for this exact version and shown on the abstract page https://arxiv.org/abs/2203.15940v3. Full licence notice: \texttt{licenses/arxiv-2203.15940-CC-BY-4.0.txt}.\par\smallskip

\noindent\textit{Editorial scope.} Counted unit: the article's theorem poteF (source0.tex lines 433--444). The article's complete original proof of it is delivered with it (source0.tex lines 445--506, one \texttt{proof} environment of 62 lines). No mathematics is written, repaired or re-derived: the statement and the proof are the article's own lines, in the article's own notation, and no retained line is substituted or re-flowed.  Retained uncounted as the source-local context the proof needs: lines 337--339; lines 341--343; lines 346--348.  Nothing was omitted from any retained span, so the package carries no omission record.  No retained line contains a citation, so the wrapper carries no bibliography and no bibliography entry.  No retained line contains a figure environment, an image-inclusion command or drawing code, so no image is delivered and the package declares no asset dependency.  Editorial formatting only, in addition: an AMS \texttt{amsart} preamble that restates the article's own declarations and macros used by the retained lines, namely the environments \texttt{lemma}, \texttt{theorem} on separate counters, together with the standard packages those lines need.  The retained environments print in this document as Lemma 1, Theorem 1 in order of appearance, whereas the article prints the same items with the numbering of its own full document.  The complete native source of the article as distributed in the arXiv e-print for this exact version is bundled unchanged as \texttt{sources/123875/source0.tex}.
\begin{equation}\label{TR}
     y(t)-y(t-r) = \int_{0}^{t} (x(s)-\Bar{x} ) ds - \int_{0}^{t-r} (x(s)-\Bar{x} ) ds = \int_{t-r}^{t} (x(s)-\Bar{x} ) ds.
\end{equation}

\begin{align}\label{detr}
    \mathcal{F}(r)&:=\sqrt{\frac{1}{(M-m)}\int _{m}^{M} \left(  y(t)-y(t-r) \right)^2  dt}.
\end{align}

\begin{lemma}\label{detrendcont}
Assume that $x\colon[0,M]\to\R$ is a continuous function. Let $0<m<M$ and consider the detrended  function $\mathcal{F}\colon[0,m]\to\R$ defined in \eqref{detr}. Then $\mathcal F^2$ is Lipschitz. In particular, $\mathcal F$ is  continuous.
\end{lemma}

\begin{theorem}\label{poteF}
    Let $M>0$ and let $x\colon [0,M]\rightarrow \R$ be a continuous function.  Let $0<m<M$ and consider the detrended function $\mathcal{F}:[0,m]\rightarrow \R$ given by \eqref{detr}.
Then $\mathcal{F}(r)$ is continuous and    
\begin{equation*}
   \lim_{r\to 0^+} \mathcal{F}(r)=0.
\end{equation*}
Moreover
\begin{equation}\label{asym}
    \lim_{r\to 0^+} \frac{\mathcal{F}(r)}{r} =\sqrt{\frac{1}{M-m}\int_{m}^{M}\left(x\left(t\right)-\overline{x}\right)^2\, dt}.
\end{equation}
In particular, $\lim_{r\to 0^+} \frac{\mathcal{F}(r)}{r}$ is a nonnegative constant number and equals 0 if and only if $x(t)=\Bar{x}$ for all $t\in[0,M]$.
\end{theorem}

\begin{proof}
By Lemma~\ref{detrendcont} the function $\mathcal F$ is right-continuous at $r=0$. Since $\mathcal F(0)=0$, we conclude that 
\begin{equation*}
    \lim_{r\to0^{+}}\mathcal{F}(r)=\mathcal{F}(0)=0.
\end{equation*}
To establish \eqref{asym}, it will suffice to prove the limit for an arbitrary sequence of positive numbers converging to zero. For this purpose, let $\{r_n\}_{n\geq 0}$ be a sequence of positive real numbers such that $r_n\to 0$ as $n\to\infty$. 
\newline
First, observe that by using \eqref{TR} we get
\[
\frac{\mathcal{F}(r)^2}{r^2}=\frac{1}{M-m}\int_{m}^{M}\biggl( \frac{1}{r}\int_{t-r}^{t}\left(x(s)-\overline{x}\right)\, ds\biggr)^2dt.
\]
Let 
\begin{equation*}
  f_n(t)=\frac{1}{r_n}\int_{t-r_n}^{t} \left(x(s)-\Bar{x}\right)ds.
\end{equation*}
Then, for a fixed $t\in[0,M]$
\begin{equation*}
    f_n(t)- \left( x(t)-\Bar{x}\right)=\frac{1}{r_n}\int_{t-r_n}^{t}\left(x(s)-x(t)\right)ds.
\end{equation*}
Since $x$ is continuous on a compact interval, it is uniformly continuous.
Let $\epsilon>0$. Then there exists $\delta>0$ such that $|s-t|<\delta$ implies $|x(s)-x(t)|<\epsilon$ for every $s$. Choose $N>0$ such that $r_n<\delta$ for every $n>N$. Then for $n>N$ and for every $s\in [t-r_n,t]$ it follows that $|x(s)-x(t)|<\epsilon.$ Hence
\begin{equation*}
   \left|\frac{1}{r_n}\int_{t-r_n}^{t}\left(x(s)-x(t)\right)ds\right|\leq \frac{1}{r_n}\int_{t-r_n}^{t}|x(s)-x(t)|ds\leq\frac{1}{r_n}\int_{t-r_n}^{t}\epsilon\, ds=\epsilon.
\end{equation*}
This proves 
\begin{equation*}
   \lim_{n\to \infty} \bigl(f_n(t)- \left( x(t)-\Bar{x}\right)\bigr)=\lim_{n\to\infty}\frac{1}{r_n}\int_{t-r_n}^{t}\left(x(s)-x(t)\right)ds=0.
\end{equation*}
Thus
\begin{equation*}
\lim_{n\to\infty}f_n(t)=x(t)-\Bar{x}.
\end{equation*}
and so
\begin{equation*}
f_n(t)^2\to\left(x(t)-\Bar{x}\right)^2  
\end{equation*}
pointwise on $[m,M].$

Using again that $x(s)$ is continuous on a compact interval we find that it is bounded. Therefore, there exists $C>0$ such that for every $s\in[0,M]$,
\begin{equation*}
    |x(s)-\Bar{x}|<C. 
\end{equation*}
We obtain
\begin{equation*}
    |f_n(t)|\leq\frac{1}{r_n}\int_{t-r_n}^{t}|x(s)-\Bar{x}|\,ds\leq\frac{1}{r_n}\int_{t-r_n}^{t}C\,ds=C.
\end{equation*}
It follows that the sequence of functions $\{f_n^2\}_{n\geq 0}$ converges pointwise to the integrable function $\left(x(t)-\Bar{x}\right)^2$ and is dominated by an
integrable (constant) function. By Lebesgue's dominated convergence theorem, we obtain
\begin{equation*}   \lim_{n\to\infty}\int_{m}^{M}f_n^2(t)\,dt=\int_{m}^{M}\lim_{n\to\infty}f_{n}^{2}(t)\, dt=\int_{m}^{M}\left(x(t)-\Bar{x}\right)^2\, dt
\end{equation*}
In summary,
\begin{align*}
   \biggl(\lim_{r\to 0^{+}}  \frac{\mathcal{F}(r)}{r}\biggr)^2&= \lim_{r\to 0^{+}}  \frac{\mathcal{F}(r)^2}{r^2} \\
   &=\frac{1}{M-m}\ \lim_{r\to 0^{+}}\int_{m}^{M}\biggl( \frac{1}{r}\int_{t-r}^{t}\left(x(s)-\overline{x}\right)\, ds\biggr)^2dt\\
   &=\frac{1}{M-m}\int_{m}^{M}\int_{m}^{M}\left(x(t)-\Bar{x}\right)^2\, dt.
\end{align*}
Hence
\begin{equation*}
    \lim_{r\to 0^+} \frac{\mathcal{F}(r)}{r} =\sqrt{\frac{1}{M-m}\int_{m}^{M}\left(x\left(t\right)-\overline{x}\right)^2\, dt},
\end{equation*} 
 as wanted.  
 \end{proof}

\end{document}
